% !TEX root = main.tex
\subsection{Instantiation Overview}
\label{subsection:overview-lattice-sms-description}

\noindent\textbf{Primitive choices.}
The generic construction is instantiated with
LP PKE~\cite{RSA:LinPei11} ($\EncryptionSchemeLP$) as $\EncryptionScheme$,
whose simple structure based on the learning-with-errors assumption is compatible with ZKBoo's arithmetic circuit proofs;
JRS signatures~\cite{CRYPTO:JeuRouSan23} ($\SignatureSchemeStandardJRS$) as~$\SignatureScheme$,
whose lattice-based signing keys support efficient zero-knowledge membership proofs;
and ZKBoo~\cite{USENIX:GiaMadOrl16} as the NIZKAoK,
with all computations over $\ZZq$ for a prime $q$.
KTX commitments~\cite{AC:KawTanXag08} replace the per-user key pair.
%when we describe membership key below.

\medskip
\noindent\textbf{OTS elimination.}
Because ZKBoo satisfies simulation-sound extractability with quasi-unique
response~\cite{INDOCRYPT:FKMV12},
it subsumes the anti-malleability role of the one-time signature $\OTSScheme$
required in the generic construction:
any adversary modifying a ZKBoo proof cannot produce a second valid proof
for a different statement without breaking the quasi-unique response property.
Retaining $\OTSScheme$ in the lattice setting would require a fresh key pair
and an additional signature over all components per signing,
increasing both signature size and signing cost;
ZKBoo's quasi-unique response gives the same guarantee at no extra cost.
Consequently, $(\vkattr, \skattr)$ no longer serves as a one-time signing key pair
for $\ovk$, but is still required for membership proofs.

\medskip
\noindent\textbf{Membership key via KTX.}
Following the KTX commitment 
\ifcameraready 
paradigm~\cite{AC:KawTanXag08},
\else 
paradigm~\cite{AC:KawTanXag08} (recalled in 
Appendix~\ref{subappendix:ktx-commitment}),
\fi
we enforce the relationship
$\vkattr = \bin(\bH \cdot \skattr)$,
where $\bH \in \ZZq^{n \times m}$ is public and random,
$\bin(\cdot)$ denotes binary decomposition,
and $\skattr$ is a uniformly random binary vector chosen by the enrolling user.
By the security of the KTX commitment, $\vkattr$ computationally hides $\skattr$;
by the hardness of the short integer solution problem, no adversary seeing $\vkattr$ can recover a binary preimage.

\medskip
\noindent\textbf{Relation $\RelationLatticeSMS$.}
The NIZKAoK proves membership in a relation $\RelationLatticeSMS$
containing two classes of constraints.

\begin{enumerate}
\item \emph{Lattice-based constraints.}
These enforce binary-vector structure, small $\ell_\infty$-norm bounds,
and linear relationships over $\ZZq$.
Binary-vector structure is captured by $\innerproduct{\bv}{\bone - \bv} = \bzero$.
Small-norm bounds are handled via the following binary decomposition lemma.

\begin{lemma}[Binary decomposition,
	{\cite{PKC:LNSW13}}, {\cite[Appendix~A.1]{ACNS:HLTW25}}]%
	\label{lemma:binary-basis-for-value-v}
Let $v \in \NN$ and $k = \lfloor \log_2 v \rfloor + 1$.
Define the basis
$\BinaryBasis_v = \{B_i\}_{i = 1}^k = \{2^i\}_{i = 0}^{k - 1} \cup \{v - (2^{k - 1} - 1)\}$.
Then $x \in [0, v]$ if and only if there exist bits $b_1, \ldots, b_k$
satisfying $x = \sum_{i = 1}^k b_i B_i$.
\end{lemma}

Using Lemma~\ref{lemma:binary-basis-for-value-v},
each bound $\|\ba\|_\infty \leq u$ is reduced to binary-vector
and linear constraints,
compatible with ZKBoo's arithmetic circuit satisfiability.
The proof of Lemma~\ref{lemma:binary-basis-for-value-v} is deferred to
\ifcameraready the full version.
\else Appendix~\ref{subappendix:preliminaries-proof-of-auxiliary-lemma}.
\fi

\item \emph{Circuit satisfiability constraints.}
These enforce $P(F(\cdot), G(\cdot))$
and any auxiliary public circuits are
correctly evaluated,
including the circuit checking $\op \in \OpeningSetSignable$.
The $P(\cdot, \cdot)$ part is handled by the PLONK reduction
of Section~\ref{section:proving-function-evaluations-via-plonk}.
\end{enumerate}

The full relation $\RelationLatticeSMS$ is 
\ifcameraready 
also deferred to the full version.
\else 
deferred to
\eqref{equation:relation-l-sms} in Appendix~\ref{subappendix:lattice-sms-description}.
\fi

\medskip
\noindent\textbf{ZKBoo as the NIZKAoK.}
ZKBoo proves arithmetic circuit satisfiability and is plausibly post-quantum.
Proving $\RelationLatticeSMS$ requires the hidden-circuit strategy of
Section~\ref{section:proving-function-evaluations-via-plonk},
which introduces additional challenges $\alpha, \beta \uniformly \ZZq$
that must be generated from the random oracle after the prover commits to $\witpre$.
The Fiat-Shamir transform~\cite{CRYPTO:FiaSha86} makes the protocol non-interactive;
however, the standard ZKBoo decomposition must be modified to accommodate
the two-phase witness structure $(\witpre, \witsuf)$.
Section~\ref{subsection:overview-lattice-sms-modifying-zkboo} describes this modification.
The result is a plausibly post-quantum \SMS scheme enforcing hidden arithmetic policies with per-signature modal accountability.